\chapter{Introduction}
\label{ch:intro}

\section{Motivation}

Mobile cellular networks must seamlessly migrate active User Equipment (UE) connections between base
stations (gNBs) without data interruption. In legacy 5G networks, handover is governed by a layer-3
Radio Resource Control (RRC) reconfiguration exchange~\cite{3gpp38300}. During this procedure, the UE
applies target configurations, releases source connectivity, and executes a contention-free Random
Access Channel (RACH) handshake on the target cell to acquire uplink timing. This process introduces
tens of milliseconds of radio interruption. While acceptable for legacy streaming, such latencies
violate the strict sub-10\,ms deadlines required by mission-critical 5G-Advanced and 6G applications,
such as extended reality (XR), vehicle-to-everything (V2X), and industrial automation.

\section{The Lower-Layer Triggered Mobility (LTM) Architecture}

To eliminate this handover delay, 3GPP Releases~18 and 19 introduced Lower-Layer Triggered
Mobility (LTM). LTM replaces the slow RRC exchange with a single layer-2 MAC Control Element (CE), the
LTM Cell Switch Command~\cite{3gpp38321}. This speed-up relies on three key mechanisms:
\begin{itemize}\itemsep2pt
  \item \textbf{Pre-cached candidate configurations:} The network provisions full configurations for
        up to eight candidate cells in advance via RRC. The switch command contains only a 3-bit index
        identifying the selected target cell.
  \item \textbf{Early timing advance acquisition:} The serving cell orders the UE to transmit a PRACH
        preamble to candidate cells before any handover. The candidate measures the preamble,
        derives a candidate Timing Advance (\TA{}), and sends it to the serving cell via F1AP to be
        cached.
  \item \textbf{RACH-less cell switch:} When the switch command contains this pre-acquired \TA{}, the
        UE applies it immediately and skips the RACH handshake entirely upon arrival at the target cell.
\end{itemize}
LTM is already undergoing deployment in commercial networks; field trials by KDDI and AT\&T have
demonstrated up to a 25\% reduction in mobility interruption~\cite{ericsson2025ltmreview,ericsson2026kddi,ericsson2026att}.

\section{Problem Statement and Research Gap}

While MAC CEs are known to lack cryptographic integrity protection (studied in 3GPP TR
33.801-01~\cite{3gpp33801} and \cite{hatinen2026overshadowing}), this report investigates a more
fundamental architectural question: what security and availability vulnerabilities are created
by the LTM architectural speed-ups themselves, independent of message forgery?

Vulnerabilities that require message forgery will become obsolete once 3GPP mandates MAC CE integrity
protection. In contrast, weaknesses stemming from unauthenticated cached state and disabled recovery
mechanisms survive any cryptographic hardening. Because LTM forms the architectural foundation for
6G mobility (whose Release~21 Stage-1 requirements begin freezing in 2027~\cite{3gpp2026rel21}, with
Stage-2/3 work continuing beyond), identifying these protocol flaws is of critical importance.

\section{The Two Findings and Methodological Approach}

By conducting a systematic, line-by-line audit of 3GPP Technical Specifications, this report uncovers
two major architectural flaws:
\begin{itemize}\itemsep2pt
  \item \textbf{Finding 1 --- Corrupted Early Timing Advance (\cref{ch:m1}):} Candidate \TA{} acquisition
        binds measurements only to a preamble index, occasion, and cell ID---never to a UE identity.
        Furthermore, 3GPP specifications impose no retention or expiration bounds on this network-held
        value. An adversary using an RF relay can forward the preamble with excess propagation delay,
        inflating the stored \TA{}. Using srsRAN and OpenAirInterface, we demonstrate that a timing error
        of 6--8 units (median 7) crosses the measured decodability boundary---deterministic under flat
        fading/AWGN, slightly more gradual under multipath---and that the relay requires only a
        $+9.9$\,dB power margin under flat fading ($+5.3$\,dB under multipath) to succeed.
  \item \textbf{Finding 2 --- Scheduling Request Deadlock (\cref{ch:m2}):} When a UE arrives via a
        RACH-less switch without a grant, it must issue Scheduling Requests (SR) on PUCCH. If these fail,
        TS~38.321 clause~5.4.4 deallocates PUCCH, SRS, and all grants unconditionally upon counter
        exhaustion. However, the subsequent instruction to recover via random access is gated by a guard
        condition that evaluates to false during an ongoing RACH-less switch. The UE is left completely
        silenced, unable to transmit or request resources.
\end{itemize}

\paragraph{The Unified Attack Composition:} Crucially, Finding~1 directly triggers Finding~2. When the
UE applies the corrupted \TA{} during a RACH-less switch, its uplink transmissions are misaligned by
6--8 units or more. Consequently, its PUCCH Format~0 Scheduling Requests cannot be decoded by the target
gNB. The SR counter exhausts with near certainty under flat fading ($P \ge 0.996$) and remains high under
multipath ($P = 0.75$--$0.97$), driving the UE into the deadlock state. Because the
deadlock blocks random access, the UE cannot re-estimate its timing on its own initiative, and a
network-issued grant fares no better since it inherits the same corrupted \TA{}---so the failure holds
until the supervision timer T304 expires.

\section{Organisation of the Report}

The remainder of this report is organized as follows: \Cref{ch:lit} reviews 5G protocol foundations, LTM
mechanics, and related literature. \Cref{ch:m1} presents Finding~1, detailing the early \TA{} vulnerability,
relay feasibility, and decodability limits. \Cref{ch:m2} presents Finding~2, formalizing the TS~38.321
scheduling request deadlock and the composite attack. \Cref{ch:results} provides experimental simulation
results, configuration sweeps, and evaluation of existing defenses. \Cref{ch:conclusion} concludes and
outlines the Phase~II research roadmap.
